\subsubsection{\texttt{fnbrk}}

\texttt{fnbrk}用来得到样条函数的相关信息

\textbf{I.}\texttt{fnbrk}的用法1为\texttt{[out1,...,outn] = fnbrk(f,part1,...,partm)},要求\texttt{n}$\leq$\texttt{m}

\begin{itemize}
    \item \texttt{parti} 表示第\texttt{i}个输入的样条信息，详见下文
    \item \texttt{f} 表示输入的样条函数
    \item \texttt{outi} 输出第\texttt{i}个输入的样条信息
\end{itemize}

对于任意格式的样条曲线 \texttt{f} ，\texttt{parti} 可以是如下形式：

\begin{table}[H]
    \centering
    \begin{tabular}{l|l}
        \hline
        \lstinline|"form"|
         & \makecell[l]{\texttt{f}的格式，返回\texttt{'ppform'}或\texttt{'B-form'}}\\
        \hline
        \lstinline|"variables"|
         & \makecell[l]{\texttt{f}的自变量的维数，返回一个整型标量}\\
        \hline
        \lstinline|"dimension"|
         & \makecell[l]{\texttt{f}的维数,返回一个整形标量}\\
        \hline
        \lstinline|"coefficients"|
         & \makecell[l]{\texttt{f}的系数,对\texttt{pp}格式，返回一个矩阵，其列数是样条函数的阶数，\\
         行数 \texttt{r = (length(breaks) - 1)}，每行是一个多项式系数的降幂排列。\\
         对 \texttt{B}格式，返回一个一维向量，其每个分量表示对应的 \texttt{B} 样条的系数}\\
        \hline
        \lstinline|"interval"|
         & \makecell[l]{\texttt{f} 的有效区间,返回一个 \texttt{1$\times$2} 的\texttt{double} 型矩阵}\\
        \hline
        \lstinline|"order"|
         & \makecell[l]{\texttt{f}的阶数,返回一个整形标量}\\
        \hline
    \end{tabular}
\end{table}

对于\texttt{pp}格式的样条曲线 \texttt{f} ，\texttt{parti} 还可以是如下形式：

\begin{table}[H]
    \centering
    \begin{tabular}{l|l}
        \hline
        \lstinline|"breaks"|
         & \makecell[l]{\texttt{f}的插值节点，返回一个一维向量}\\
        \hline
        \lstinline|"pieces"|
         & \makecell[l]{\texttt{f}的多项式段数，返回一个整型标量}\\
        \hline
    \end{tabular}
\end{table}

对于\texttt{B}格式的样条曲线f ，\texttt{parti} 还可以是如下形式：

\begin{table}[H]
    \centering
    \begin{tabular}{l|l}
        \hline
        \lstinline|"knots"|
         & \makecell[l]{\texttt{f}的\texttt{B}样条节点，返回一个一维向量}\\
        \hline
        \lstinline|"number"|
         & \makecell[l]{\texttt{f}的系数的数量，返回一个整型标量}\\
        \hline
    \end{tabular}
\end{table}

\textbf{I.}
输入：
\begin{lstlisting}[language = matlab]
    x = [  0,2,3, 5,6  ];
    y = [1,0,3,1,-1,0,0];
    f = csape(x,y,'complete');
    [form,variables,dimension,coefs,interval,order,breaks,pieces]...
    =fnbrk(f,'form','variables','dimension','coefficients','interval','order','breaks','pieces')
\end{lstlisting}

 输出为：
\begin{lstlisting}[language = matlab ]
         form = 'ppform'
    variables = 1
    dimension = 1
        coefs = 
                -0.6499    1.5497    1.0000         0
                 0.9489   -2.3495   -0.5995    3.0000
                 0.1142    0.4973   -2.4516    1.0000
                -1.0914    1.1828    0.9086   -1.0000
     interval = 0   6
        order = 4
       breaks = 0   2   3   5   6
       pieces = 4
\end{lstlisting}

\textbf{II.}
输入：
\begin{lstlisting}[language = matlab]
knots = [0,0,0,0,1,2,2,2,2];
    x = [0,1,1,1,2];
    y = [2,0,1,2,-1];
    f = spapi(knots,x,y);
    [form,variables,dimension,coefs,interval,order,knots,number] = ...
    fnbrk(f,'form','variables','dimension','coefficients','interval','order','knots','number')
\end{lstlisting}

 输出为：
\begin{lstlisting}[language = matlab ]
         form = 'B-form'
    variables = 1
    dimension = 1
        coefs = 
                2.0000   -0.3333   -0.3333    1.0000   -1.0000
     interval = 0   2
        order = 4
        knots = 0   0   0   0   1   2   2   2   2
       number = 5
\end{lstlisting}

\textbf{II.}
\texttt{fnbrk}的用法2为\texttt{out=fnbrk(f,interval)}，可以与用法1混用，即\texttt{f} 之后的参数类型既可以是矩阵，也可以是字符串
\begin{itemize}
    \item \texttt{interval} 为一个大小为1 $\times $2的矩阵，表示区间
    \item \texttt{f} 表示输入的样条函数
    \item \texttt{out} 输出\texttt{f}截取在区间\texttt{interval}的\texttt{pp}格式样条结构体
\end{itemize} 
